\begin{python}
from scripts.calc_4 import *
\end{python}

\section{Составление уравнений состояния цепи}

Для составления уравнений состояния выполним замену индуктивности источниками тока $i_{L1}$ и $i_{L2}$, а емкость --- источником напряжения $u_{C1}$.

\begin{figure}[H]
    \centering
    \begin{tikzpicture}[transform shape]
        % Paths, nodes and wires:
        \draw (1, 1) to[american current source, l={$i_1$}] (1, 4);
        \draw (3, 4) to[european resistor, l_={$R_1$}] (3, 1);
        \draw (8, 4) to[european resistor, l={$R_2$}] (8, 1);
        \draw (1, 1) -- (3, 1) -- (5, 1) -- (8, 1);
        \draw (1, 4) -- (3, 4) -- (5, 4);
        \node[inputarrow, rotate=-90] at (3, 3.25){};
        \node[inputarrow, rotate=-90] at (8, 3.25){};
        \draw (5, 4) to[american current source, l={$i_{L2}$}] (8, 4);
        \draw (5, 4) to[american current source, l={$i_{L1}$}] (5, 2.5);
        \draw (5, 2.5) to[american voltage source, l={$u_{C1}$}] (5, 1);
        \node[shape=rectangle, minimum width=0.59cm, minimum height=0.465cm](N1) at (3.5, 3){} node[anchor=center] at (N1.text){$+$} node[anchor=north west, align=left, text width=0.202cm, inner sep=6pt] at (3.188, 3.25){};
        \node[shape=rectangle, minimum width=0.59cm, minimum height=0.465cm](N2) at (3.5, 2){} node[anchor=center] at (N2.text){$-$} node[anchor=north west, align=left, text width=0.202cm, inner sep=6pt] at (3.188, 2.25){};
        \node[shape=rectangle, minimum width=0.59cm, minimum height=0.465cm](N3) at (7.5, 3){} node[anchor=center] at (N3.text){$+$} node[anchor=north west, align=left, text width=0.202cm, inner sep=6pt] at (7.188, 3.25){};
        \node[shape=rectangle, minimum width=0.59cm, minimum height=0.465cm](N4) at (7.5, 2){} node[anchor=center] at (N4.text){$-$} node[anchor=north west, align=left, text width=0.202cm, inner sep=6pt] at (7.188, 2.25){};
    \end{tikzpicture}
    \caption{Схема замещения для составления уравнений состояния}
    \label{fig:scheme_state}
\end{figure}

Запишем систему уравнений состояния:
\begin{equation}
    \begin{cases}
        u_{C1}' = \frac{1}{C_1} i_{C1} \\
        i_{L1}' = \frac{1}{L_1} u_{L1} \\
        i_{L2}' = \frac{1}{L_2} u_{L2} \\
    \end{cases}
    \label{form:уравнение_состояния}
\end{equation}

Ток через конденсатор:
\begin{equation}
    i_{C1} = i_{L1}
    \label{form:ic1}
\end{equation}

Напряжение на первой индуктивности определяется через контур, включающий $R_1$ и $C_1$:
\begin{equation}
    u_{L1} = u_{R1} - u_{C1}
\end{equation}

Напряжение на резисторе $R_1$ зависит от узловых токов:
\begin{equation}
    u_{R1} = (i_1 - i_{L1} - i_{L2})R_1
\end{equation}

Следовательно:
\begin{equation}
    u_{L1} = (i_1 - i_{L1} - i_{L2})R_1 - u_{C1}
    \label{form:ul1}
\end{equation}

Напряжение на второй индуктивности:
\begin{equation}
    u_{L2} = u_{R1} - u_{R2} = (i_1 - i_{L1} - i_{L2})R_1 - i_{L2}R_2
    \label{form:ul2}
\end{equation}

Зная \cref{form:ic1,form:ul1,form:ul2} сформируем уравнения состояния по образцу \cref{form:уравнение_состояния}:
\begin{equation}
    \begin{cases}
        u_{C1}' = \frac{1}{C_1} i_{L1} \\
        i_{L1}' = -\frac{1}{L_1} u_{C1} - \frac{R_1}{L_1} i_{L1} - \frac{R_1}{L_1} i_{L2} + \frac{R_1}{L_1} i_1 \\
        i_{L2}' = -\frac{R_1}{L_2} i_{L1} - \frac{R_1 + R_2}{L_2} i_{L2} + \frac{R_1}{L_2} i_1
    \end{cases}
    \label{form:uravn_sost}
\end{equation}

Подставляя нормированные значения параметров, система уравнений состояния записывается в матричной форме:

\begin{equation}
    \begin{bmatrix}
        u_{C1}' \\
        i_{L1}' \\
        i_{L2}'
    \end{bmatrix}
    =
    \begin{bmatrix}
        \pv[.2]{mat['A11']} & \pv[.2]{mat['A12']} & \pv[.2]{mat['A13']} \\
        \pv[.2]{mat['A21']} & \pv[.2]{mat['A22']} & \pv[.2]{mat['A23']} \\
        \pv[.2]{mat['A31']} & \pv[.2]{mat['A32']} & \pv[.2]{mat['A33']}
    \end{bmatrix}
    \cdot
    \begin{bmatrix}
        u_{C1} \\
        i_{L1} \\
        i_{L2}
    \end{bmatrix}
    +
    \begin{bmatrix}
        \pv[.2]{mat['B11']} \\
        \pv[.2]{mat['B21']} \\
        \pv[.2]{mat['B31']}
    \end{bmatrix}
    \cdot
    i_1
    \label{form:uravn_sost_matrix}
\end{equation}